\section{Einleitung}
\label{sec:introduction}

\subsection{Inhaltliche Stichpunkte}
\begin{itemize}
  \item \textbf{Motivation}
  \begin{itemize}
    \item Zunehmende Vernetzung von Ladeinfrastruktur für Elektrofahrzeuge
          (EVSE) und daraus entstehende Cybersicherheitsrisiken.
    \item Im Datensatz werden acht Knoten (Node A bis H) und die Klassen
          Backdoor, syn-flood und ``none'' untersucht.
  \end{itemize}
  \item \textbf{Problemstellung}
  \begin{itemize}
    \item Lokale Beobachtungen eines einzelnen Knotens können Hinweise auf
          koordinierte oder über mehrere Knoten verteilte Angriffe verpassen.
    \item Forschungsbedarf und Aussagen über bestehende IDS mit Literatur
          belegen; diese Begründung ist nicht durch die Evaluationsnotebooks
          allein abgedeckt.
  \end{itemize}
  \item \textbf{Lösungsansatz}
  \begin{itemize}
    \item Neighbor Context Detection: lokale Messwerte und
          Erkennungswahrscheinlichkeiten werden durch aggregierte
          Informationen anderer Knoten ergänzt.
    \item Im ausgewerteten Aggregationslauf gehen die Informationen aller
          sieben übrigen Knoten ein; die genaue Topologie ist in
          Kapitel~\ref{sec:methodology} zu klären.
  \end{itemize}
  \item \textbf{Forschungsfrage}
  \begin{itemize}
    \item Führt der Austausch von Kontextinformationen benachbarter Knoten
          in einem P2P-Netzwerk zu einer signifikant besseren
          Erkennungsleistung?
    \item Die Notebook-Ergebnisse zeigen höhere Kennzahlen für das
          Full-Context-Modell; ein statistischer Signifikanztest ist dort
          jedoch nicht dokumentiert.
  \end{itemize}
  \item \textbf{Beitrag der Arbeit}
  \begin{itemize}
    \item Evaluation einer lokalen Baseline gegenüber Random-Forest-Modellen
          mit aggregiertem Nachbarschaftskontext.
    \item Hybride Datengenerierung mit realen Daten sowie klassenspezifischen
          Anteilen von TabGAN- und DAE-Daten.
    \item Evaluation über acht simulierte Knoten und mehrere Angriffsszenarien.
  \end{itemize}
\end{itemize}

\subsection{Noch zu erklären}
\begin{itemize}
  \item Was ``signifikant'' in der Forschungsfrage bedeutet und welche
        statistischen Tests oder Wiederholungen dazu noch durchgeführt werden.
  \item Ob die untersuchten Eingaben tatsächlich Netzwerkverkehr oder
        elektrische Betriebsmerkmale der Ladepunkte darstellen: Die
        Evaluationsnotebooks verwenden shunt\_voltage, bus\_voltage\_V,
        current\_mA, power\_mW und State.
  \item Welche Aussagen durch eigene Messungen belegt werden und welche
        durch Literatur begründet werden müssen.
\end{itemize}

\subsection{Grafikbedarf}
\begin{itemize}
  \item \textbf{Übersichtsdiagramm:} acht EVSE-Knoten, lokale Erkennung,
        Austausch aggregierter Kontextinformationen und die drei Zielklassen.
\end{itemize}
